\section{Exact controls and proves-too-much boundary}
\label{sec:controls}

The main statements are analytic identities, so their decisive controls can be
proved without a root-fitting experiment.  Let $1<x_1<x_2<\cdots$ be any
increasing inventory and replace $(p_k)$ by $(x_k)$.  Adjacent pairing still
gives
\[
 \sum_n|x_{2n-1}^{-s}-x_{2n}^{-s}|
 \le |s|\int_{x_1}^{\infty}x^{-\RePart s-1}\dd x.                  \tag{7.1}
\]
Thus the half-plane theorem, relative product, and zero-free reflected strip
hold for sorted composites, random integers, perturbed entropies, and matched
non-prime inventories.  Only the interpretation of the labels changes.

\begin{center}\small
\begin{tabularx}{0.98\textwidth}{L{0.24\textwidth}L{0.24\textwidth}X}
\toprule
Control & Analytic result & Arithmetic diagnosis\\\midrule
Reverse $+|-$ orientation & $R\mapsto R^{-1}$; same half-plane & no intrinsic choice of positive Euler orientation\\
Shift pairing by one rank & same trace-class estimate, up to finite factors & most signs change although the inventory does not\\
Balanced $m$-blocks & full half-plane iff block sum is zero & every admissible fixed pattern contains cancellation\\
Sorted composites & identical interval proof & no prime-specific discrimination\\
Random matched inventory & identical interval proof after sorting & mechanism proves too much\\
All-positive blocks & fails at zero order & correct orientation loses the common strip\\
Unitary phase $u_p$ & repetition sign becomes $u_p^r$ & different candidate and determinant\\
\bottomrule
\end{tabularx}
\end{center}

The two-atom prefix already displays every structural effect.  For $a<b$,
\[
 R_{a,b}(s,z)=\frac{1-za^{-s}}{1-zb^{-s}},
 \qquad
 H_{a,b}(s,z)=R_{a,b}(s,z)R_{a,b}(1-s,z).       \tag{7.2}
\]
At $z=1$ and $0<\RePart s<1$, all four scalar factors are nonzero.  On the
critical line the modulus generally moves, but there is no zero.  Adding more
adjacent pairs preserves these properties through a locally uniform relative
product.

Shifted pairing is particularly adversarial.  Pairing
$(p_2,p_3),(p_4,p_5),\ldots$ obeys the same proof, while the unpaired factor
at $p_1$ is only a finite-rank correction.  The parity of almost every prime
is reversed.  Hence the analytic continuation mechanism cannot decide which
prime factors belong in the numerator.  Entropy ordering chooses a reproducible
pattern, but the trace-class theorem does not certify its arithmetic sign.

Larger blocks reach the same conclusion with more freedom.  Coefficients such
as $(1,1,-1,-1)$ or $(1,\omega,\omega^2)$, with $\omega^3=1$, satisfy zero
block sum and therefore inherit the right-half-plane estimate.  Their
prime-power ledgers are mutually incompatible.  Constant phases can change
the block pattern, but they cannot turn a zero-sum vector into the all-positive
vector.

These controls trigger the Route-A adversarial verdict
\[
 \boxed{\texttt{STOP\_SCOPED / PROVES\_TOO\_MUCH}}.               \tag{7.3}
\]
The proved theorem remains useful as a symbolic regularization principle.  It
cannot, by itself, distinguish the rational primes or infer a critical-line
divisor.
